\section{Introduction}
\label{sec:intro}

Seasonal snow is a key component of the mountain water cycle. In
Mediterranean mountain ranges such as the Pyrenees, snowmelt sustains river
discharge, reservoir storage, hydropower generation and irrigation during the
dry season, and the spatial distribution of the snowpack governs avalanche
hazard, soil moisture and the timing of vegetation growth
\citep{lopezmoreno2013,fayad2017}. Snow depth (HS) is highly variable in
space: wind redistribution, terrain shading, slope, aspect and preferential
deposition produce sharp gradients over distances of a few metres, especially
in wind-exposed alpine terrain \citep{mott2018,revuelto2014}. Capturing this
fine-scale variability is therefore essential, but it is precisely where
coarse-resolution products and point measurements fail.

Airborne and terrestrial Light Detection and Ranging (LiDAR) surveys have
become the reference technique for mapping HS at sub-metre resolution by
differencing snow-on and snow-off digital surface models
\citep{deems2013,painter2016}. However, LiDAR acquisitions are expensive and
logistically demanding, so they can only be carried out on a limited number of
dates per season. This motivates the development of predictive models that can
\emph{interpolate and extrapolate} HS from a small set of LiDAR surveys, using
spatially exhaustive and inexpensive predictors such as terrain attributes,
satellite-derived snow-cover information and meteorological forcing.

Two broad families of methods have been applied to this problem. Statistical
and machine-learning regressors (multiple linear regression, generalised
additive models, Random Forests, gradient boosting) operate on a per-pixel
basis and have proven effective at relating HS to topographic predictors
\citep{revuelto2014,broxton2019}. Their main limitation is that they treat
each pixel independently and therefore cannot exploit the spatial context of
the snowpack. Convolutional neural networks (CNNs), by contrast, learn
spatially structured representations and have recently been used for snow and
related cryospheric mapping tasks \citep{daudt2018,king2020}. Yet most studies
still evaluate model quality almost exclusively with pixel-wise metrics --
coefficient of determination (\Rsq{}), root mean squared error (RMSE), mean
absolute error (MAE) -- which reward the correct prediction of the mean and
penalise outliers, but are largely insensitive to whether the model reproduces
the \emph{spatial pattern} of accumulation and scour. A model can attain a
respectable \Rsq{} while producing an over-smoothed map that misses the very
drift features that matter for avalanche and hydrological applications.

In this work we address both the modelling and the evaluation sides of this
gap. We build a high-resolution (1\,m) HS prediction benchmark for the Izas
experimental catchment in the Central Spanish Pyrenees, using a 27-date LiDAR
time series (2021--2025) provided in collaboration with the Pyrenean
Institute of Ecology (IPE--CSIC). On top of this benchmark we (i) compare a
Random Forest baseline with two CNN architectures (U-Net and ResUNet++) under
a strict temporal train/validation/test split, (ii) adopt the Spatial
Efficiency metric (\SPAEF{}) and a modified, bias-aware variant (\MSPAEF{}) as
pattern-oriented evaluation criteria alongside the conventional pixel-wise
metrics, and (iii) introduce a hybrid spatial loss that explicitly optimises
for spatial pattern fidelity through a tunable correlation term.

\subsection{Contributions}
\label{subsec:contributions}

The main contributions of this paper are:

\begin{enumerate}
  \item \textbf{A reproducible 1\,m snow-depth benchmark} for a Pyrenean
  catchment, built from airborne LiDAR and a rich stack of 22 topographic,
  snow-cover and wind-redistribution predictors, with a temporally
  disjoint test year (2025) that stresses the extrapolation ability of the
  models.

  \item \textbf{A pattern-aware evaluation protocol} based on \SPAEF{} and
  \MSPAEF{}, which we show to be far more discriminative than \Rsq{} for this
  task -- in particular, it exposes a ranking reversal between the Random Forest
  and the U-Net (the RF leads on \Rsq{} but the U-Net leads on \SPAEF{}) and the
  strong seed-instability of the plain U-Net, both of which a single-run,
  \Rsq{}-only assessment would hide.

  \item \textbf{A hybrid spatial loss} combining pixel-wise mean squared error
  with a spatial Pearson-correlation term weighted by a hyperparameter
  $\lpear$, together with a complete sweep ($\lpear\in[0,1]$, eight values,
  three seeds) that identifies $\lpear=0.4$ as the optimum, improving both
  \Rsq{} and \SPAEF{} while halving the variance across seeds.

  \item \textbf{A leave-one-group-out ablation} of the predictor channels that
  quantifies the relative importance of fine/coarse topography, snow
  persistence, wind sheltering and satellite snow-cover extent.

  \item \textbf{An analysis of meteorological forcing} (air temperature and
  accumulated precipitation as additional channels), showing that, for this
  catchment and predictor set, meteorological inputs do not significantly
  improve the prediction once topography and snow history are available.
\end{enumerate}

\subsection{Paper organisation}
\label{subsec:organisation}

The remainder of the paper is organised as follows.
Section~\ref{sec:related} reviews related work on snow-depth modelling, deep
learning for environmental mapping and spatial evaluation metrics.
Section~\ref{sec:data} describes the study area, the LiDAR-derived dataset and
the predictor channels.
Section~\ref{sec:methods} presents the models, the \SPAEF{}/\MSPAEF{} metrics
and the spatial loss, together with the experimental protocol.
Section~\ref{sec:results} reports the model comparison, the $\lpear$ sweep, the
channel ablation and the meteorological experiments.
Section~\ref{sec:discussion} discusses the implications and limitations, and
Section~\ref{sec:conclusions} concludes.
